% GENERATED FILE. Edit canonical lessons, then run npm run generate:books.
% canonical-source-hash: fnv1a64:b54d5a0bf2a49a45
% canonical-lessons: HI-C341-panna, HI-C341-jild, HI-C341-sirsak, HI-C341-pathak, HI-C341-bikna

\chapter{Turning the Pages}
\label{ch:hi-turning-the-pages}
\hlchaptermodality{341}

\begin{chapteropening}
\textbf{By the end of this chapter:} \emph{I can describe a book and say how well it sold.}

Complete the last lesson of chapter 341: I can describe a book and say how well it sold.
\end{chapteropening}

\section[\texorpdfstring{\hi{पन्ना} (pannā)}{पन्ना (pannā)}]{\hi{पन्ना} (pannā) — a page; a leaf (of a book)}

\label{lesson:HI-C341-panna}



\begin{quote}
Before the new one: say the Hindi for a chapter of a book, then the Hindi for an outcome.
\end{quote}

\subsection*{You'll want to know: \hi{पन्ना}}
\textbf{\hi{पन्ना}} — \emph{pannā} — \textquotedblleft{}a page; a leaf (of a book)\textquotedblright{}.

\begin{quote}
\hi{बारिश} \hi{में} \hi{इस} \hi{किताब} \hi{के} \hi{कुछ} \hi{पन्ने} \hi{फट} \hi{गए।}
\end{quote}

\begin{quote}
\emph{bāriś meṁ is kitāb ke kuch panne phaṭ gae.}
\end{quote}

\begin{quote}
Some pages of this book got torn in the rain.
\end{quote}

Masculine.

\begin{cousinweb}
\textbf{\hi{पन्ना}} (\emph{pannā}) goes back to Sanskrit \emph{parṇa}, \textquotedblleft{}a leaf\textquotedblright{} (CDIAL 7918), where CDIAL gives Hindi \emph{pannā}, a leaf or a page. A page of a book is a leaf, as in English.
\end{cousinweb}

\begin{grammarlens}[title={What you've built}]
One more word for telling what happened.
\end{grammarlens}

\subsection*{Your turn}
\begin{itemize}
\raggedright
  \item \emph{Say it:} \emph{pannā}
  \item \emph{Say it:} \emph{pannā}, once more
  \item \emph{Say it:} \emph{bāriś meṁ is kitāb ke kuch panne phaṭ gae.} — Some pages of this book got torn in the rain.
  \item \emph{Recall:} say the Hindi for a chapter of a book, then the Hindi for an outcome, then say \emph{pannā} again
\end{itemize}

\begin{tcolorbox}[breakable,colback=teal!4,colframe=teal!35!black,title={Before you move on}]
What does \hi{पन्ना} mean? (\textquotedblleft{}A page; a leaf (of a book)\textquotedblright{}.) Say it once more.
\end{tcolorbox}
\section[\texorpdfstring{\hi{जिल्द} (jild)}{जिल्द (jild)}]{\hi{जिल्द} (jild) — a binding; a cover (of a book)}

\label{lesson:HI-C341-jild}



\begin{quote}
Before the new one: say the Hindi for an outcome, then the Hindi for a page.
\end{quote}

\subsection*{You'll want to know: \hi{जिल्द}}
\textbf{\hi{जिल्द}} — \emph{jild} — \textquotedblleft{}a binding; a cover (of a book)\textquotedblright{}.

\begin{quote}
\hi{पुरानी} \hi{किताब} \hi{की} \hi{जिल्द} \hi{टूट} \hi{गई।}
\end{quote}

\begin{quote}
\emph{purānī kitāb kī jild ṭūṭ gaī.}
\end{quote}

\begin{quote}
The old book's binding came apart.
\end{quote}

Feminine.

\begin{cousinweb}
\textbf{\hi{जिल्द}} (\emph{jild}) is usually traced, through Persian, to Arabic \emph{jild}, \textquotedblleft{}skin, hide\textquotedblright{}: books were once bound in leather.
\end{cousinweb}

\begin{grammarlens}[title={What you've built}]
One more word for telling what happened.
\end{grammarlens}

\subsection*{Your turn}
\begin{itemize}
\raggedright
  \item \emph{Say it:} \emph{jild}
  \item \emph{Say it:} \emph{jild}, once more
  \item \emph{Say it:} \emph{purānī kitāb kī jild ṭūṭ gaī.} — The old book's binding came apart.
  \item \emph{Recall:} say the Hindi for an outcome, then the Hindi for a page, then say \emph{jild} again
\end{itemize}

\begin{tcolorbox}[breakable,colback=teal!4,colframe=teal!35!black,title={Before you move on}]
What does \hi{जिल्द} mean? (\textquotedblleft{}A binding; a cover (of a book)\textquotedblright{}.) Say it once more.
\end{tcolorbox}
\section[\texorpdfstring{\hi{शीर्षक} (śīrṣak)}{शीर्षक (śīrṣak)}]{\hi{शीर्षक} (śīrṣak) — a title; a heading}

\label{lesson:HI-C341-sirsak}



\begin{quote}
Before the new one: say the Hindi for a page, then the Hindi for a binding.
\end{quote}

\subsection*{You'll want to know: \hi{शीर्षक}}
\textbf{\hi{शीर्षक}} — \emph{śīrṣak} — \textquotedblleft{}a title; a heading\textquotedblright{}.

\begin{quote}
\hi{मेरी} \hi{राय} \hi{में} \hi{उपन्यास} \hi{का} \hi{शीर्षक} \hi{बहुत} \hi{छोटा} \hi{है।}
\end{quote}

\begin{quote}
\emph{merī rāy meṁ upanyās kā śīrṣak bahut choṭā hai.}
\end{quote}

\begin{quote}
In my opinion the novel's title is too short.
\end{quote}

Masculine.

\begin{cousinweb}
\textbf{\hi{शीर्षक}} (\emph{śīrṣak}) is Sanskrit \emph{śīrṣaka} (Monier-Williams), from \emph{śīrṣa}, the head: what stands at the head of a book or a page.
\end{cousinweb}

\begin{grammarlens}[title={What you've built}]
One more word for telling what happened.
\end{grammarlens}

\subsection*{Your turn}
\begin{itemize}
\raggedright
  \item \emph{Say it:} \emph{śīrṣak}
  \item \emph{Say it:} \emph{śīrṣak}, once more
  \item \emph{Say it:} \emph{merī rāy meṁ upanyās kā śīrṣak bahut choṭā hai.} — In my opinion the novel's title is too short.
  \item \emph{Recall:} say the Hindi for a page, then the Hindi for a binding, then say \emph{śīrṣak} again
\end{itemize}

\begin{tcolorbox}[breakable,colback=teal!4,colframe=teal!35!black,title={Before you move on}]
What does \hi{शीर्षक} mean? (\textquotedblleft{}A title; a heading\textquotedblright{}.) Say it once more.
\end{tcolorbox}
\section[\texorpdfstring{\hi{पाठक} (pāṭhak)}{पाठक (pāṭhak)}]{\hi{पाठक} (pāṭhak) — a reader}

\label{lesson:HI-C341-pathak}



\begin{quote}
Before the new one: say the Hindi for a binding, then the Hindi for a title.
\end{quote}

\subsection*{You'll want to know: \hi{पाठक}}
\textbf{\hi{पाठक}} — \emph{pāṭhak} — \textquotedblleft{}a reader\textquotedblright{}.

\begin{quote}
\hi{उस} \hi{किताब} \hi{के} \hi{लाखों} \hi{पाठक} \hi{हैं।}
\end{quote}

\begin{quote}
\emph{us kitāb ke lākhoṁ pāṭhak haiṁ.}
\end{quote}

\begin{quote}
That book has hundreds of thousands of readers.
\end{quote}

Masculine.

\begin{cousinweb}
\textbf{\hi{पाठक}} (\emph{pāṭhak}) is Sanskrit \emph{pāṭhaka}, \textquotedblleft{}a reader, a reciter\textquotedblright{} (Monier-Williams), the same family as \textbf{\hi{पाठ}}, a lesson.
\end{cousinweb}

\begin{grammarlens}[title={What you've built}]
One more word for telling what happened.
\end{grammarlens}

\subsection*{Your turn}
\begin{itemize}
\raggedright
  \item \emph{Say it:} \emph{pāṭhak}
  \item \emph{Say it:} \emph{pāṭhak}, once more
  \item \emph{Say it:} \emph{us kitāb ke lākhoṁ pāṭhak haiṁ.} — That book has hundreds of thousands of readers.
  \item \emph{Recall:} say the Hindi for a binding, then the Hindi for a title, then say \emph{pāṭhak} again
\end{itemize}

\begin{tcolorbox}[breakable,colback=teal!4,colframe=teal!35!black,title={Before you move on}]
What does \hi{पाठक} mean? (\textquotedblleft{}A reader\textquotedblright{}.) Say it once more.
\end{tcolorbox}
\section[\texorpdfstring{\hi{बिकना} (biknā)}{बिकना (biknā)}]{\hi{बिकना} (biknā) — to be sold; to sell}

\label{lesson:HI-C341-bikna}



\begin{quote}
Before the new one: say the Hindi for a title, then the Hindi for a reader.
\end{quote}

\subsection*{You'll want to know: \hi{बिकना}}
\textbf{\hi{बिकना}} — \emph{biknā} — \textquotedblleft{}to be sold; to sell\textquotedblright{}.

\begin{quote}
\hi{पहले} \hi{हफ़्ते} \hi{में} \hi{ही} \hi{वह} \hi{उपन्यास} \hi{बहुत} \hi{बिका।}
\end{quote}

\begin{quote}
\emph{pahle hafte meṁ hī vah upanyās bahut bikā.}
\end{quote}

\begin{quote}
The novel sold very well in its very first week.
\end{quote}

Past (it …): \hi{वह} \hi{किताब} \hi{बिक} \hi{गई}, \emph{vah kitāb bik gaī}.

\begin{cousinweb}
\textbf{\hi{बिकना}} (\emph{biknā}) goes back to Sanskrit \emph{vikrīyate}, \textquotedblleft{}is sold\textquotedblright{} (CDIAL 11642). It is the intransitive partner of \textbf{\hi{बेचना}}: you \textbf{\hi{बेचते} \hi{हो}} a book, and the book \textbf{\hi{बिकती} \hi{है}}.
\end{cousinweb}

\begin{grammarlens}[title={What you've built}]
One more word for telling what happened.
\end{grammarlens}

\subsection*{Your turn}
\begin{itemize}
\raggedright
  \item \emph{Say it:} \emph{biknā}
  \item \emph{Say it:} \emph{biknā}, once more
  \item \emph{Say it:} \emph{pahle hafte meṁ hī vah upanyās bahut bikā.} — The novel sold very well in its very first week.
  \item \emph{Recall:} say the Hindi for a title, then the Hindi for a reader, then say \emph{biknā} again
\end{itemize}

\begin{tcolorbox}[breakable,colback=teal!4,colframe=teal!35!black,title={Before you move on}]
What does \hi{बिकना} mean? (\textquotedblleft{}To be sold; to sell\textquotedblright{}.) Say it once more.
\end{tcolorbox}
